\documentclass[12pt]{scrartcl}

\input{../styles/Packages}
\input{../styles/FormatAndHeader}

\usepackage{xcolor} % Für die Farben im Code

% Definiere das Aussehen des Codes
\lstset{
    language=SQL,
    basicstyle=\ttfamily\small,
    keywordstyle=\color{blue}\bfseries,
    stringstyle=\color{red},
    commentstyle=\color{gray}\itshape,
    numbers=left,           % Zeilennummern links
    numberstyle=\tiny\color{gray},
    stepnumber=1,
    showstringspaces=false, % Keine komischen Unterstriche in Strings
    tabsize=4,
    breaklines=true,        % Automatischer Zeilenumbruch
    breakatwhitespace=false,
    frame=single,           % Ein Rahmen um den Code
    extendedchars=true,
    literate={ö}{{\"o}}1 {ä}{{\"a}}1 {ü}{{\"u}}1 {ß}{{\ss}}1 {Ö}{{\"O}}1 {Ä}{{\"A}}1 {Ü}{{\"U}}1
}

% \stepcounter{countername}: Increases counter
\setcounter{sheetnr}{4} % Number of sheet


\setcounter{exnum}{1} % Exercise number

\begin{document}
% Nutze den \exercise{Aufgabenname} Befehl, um eine neue Aufgabe zu beginnen.
% Möchtest du eine Aufgabe in der Nummerierung überspringen, schreibe vor der Aufgabe: \stepcounter{exnum}
% Möchtest du die Nummer einer Aufgabe auf eine beliebige Zahl x setzen, schreibe vor der Aufgabe: \setcounter{exnum}{x}

    %
    \exercise{Abgabe 4.1: SQL-Abfragen}
    \exercise{4.2: Relationenalgebra und SQL}

\subsection*{a) Relationenalgebra: Dozenten, die keine Vorlesung halten}
% Wir ermitteln die Personalnummern aller Professoren und ziehen davon die
% Personalnummern derer ab, die in "Vorlesungen" als "gelesenVon" eingetragen sind.
% Anschließend wird das Ergebnis mit Professoren gejoint, um den Namen zu projizieren.
\[
\pi_{Name}\left( Professoren \bowtie \left( \pi_{PersNr}(Professoren) - \pi_{gelesenVon}(Vorlesungen) \right) \right)
\]

\subsection*{b) SQL: Dozenten, die keine Vorlesung halten}
% Umsetzung der Differenz aus Teilaufgabe a) mittels einer NOT IN-Unterabfrage.
\begin{lstlisting}[language=SQL]
SELECT Name
FROM Professoren
WHERE PersNr NOT IN (
    SELECT gelesenVon
    FROM Vorlesungen
    WHERE gelesenVon IS NOT NULL
);
\end{lstlisting}

\subsection*{c) Relationenalgebra: Studenten mit mindestens zwei Vorlesungen von Kant}
% Wir nutzen zwei unabhängige Kopien (mittels Umbenennung \rho) der "hoeren"-Relation,
% um paarweise unterschiedliche Vorlesungen (V1.VorlNr \neq V2.VorlNr) für dieselbe
% MatrNr zu finden, die beide von Kant gelesen werden.
\begin{align*}
&\pi_{Name}\Big( Studenten \bowtie \pi_{MatrNr} \Big( \\
&\quad  (\rho_{V1}(hoeren) \bowtie \sigma_{Name=\text{'Kant'}}(Professoren \bowtie_{PersNr=gelesenVon} Vorlesungen)) \\
&\quad  \bowtie_{V1.MatrNr = V2.MatrNr \land V1.VorlNr \neq V2.VorlNr} \\
&\quad  (\rho_{V2}(hoeren) \bowtie \sigma_{Name=\text{'Kant'}}(Professoren \bowtie_{PersNr=gelesenVon} Vorlesungen)) \\
&\Big)\Big)
\end{align*}

\subsection*{d) SQL: Studenten mit mindestens zwei Vorlesungen von Kant}
% Eleganter und effizienter über GROUP BY und HAVING COUNT.
% Es werden alle Bedingungen der Joins explizit im WHERE-Teil deklariert.
\begin{lstlisting}[language=SQL]
SELECT Studenten.Name
FROM Studenten, hoeren, Vorlesungen, Professoren
WHERE Studenten.MatrNr = hoeren.MatrNr
  AND hoeren.VorlNr = Vorlesungen.VorlNr
  AND Vorlesungen.gelesenVon = Professoren.PersNr
  AND Professoren.Name = 'Kant'
GROUP BY Studenten.MatrNr, Studenten.Name
HAVING COUNT(hoeren.VorlNr) >= 2;
\end{lstlisting}

\subsection*{e) Relationenalgebra: Studenten, die nur Vorlesungen mit 4 SWS hören}
% All-Quantifizierung ("nur") wird über die Mengendifferenz realisiert:
% Alle Studenten, die überhaupt irgendetwas hören minus diejenigen Studenten,
% die mindestens eine Vorlesung hören, die NICHT 4 SWS hat.
\[
\pi_{Name}\left( Studenten \bowtie \left(
    \pi_{MatrNr}(hoeren) - \pi_{MatrNr}(\sigma_{SWS \neq 4}(hoeren \bowtie Vorlesungen))
\right)\right)
\]

\subsection*{f) SQL: Studenten, die nur Vorlesungen mit 4 SWS hören}
% Analoge Umsetzung zur Relationenalgebra mittels NOT IN und einer Unterabfrage,
% die alle Hörer von Vorlesungen ungleich 4 SWS identifiziert.
\begin{lstlisting}[language=SQL]
SELECT Name
FROM Studenten
WHERE MatrNr IN (SELECT MatrNr FROM hoeren) -- Hat ueberhaupt Vorlesungen
  AND MatrNr NOT IN (
    SELECT hoeren.MatrNr
    FROM hoeren, Vorlesungen
    WHERE hoeren.VorlNr = Vorlesungen.VorlNr
      AND Vorlesungen.SWS <> 4
  );
\end{lstlisting}

    \newpage

    \exercise{Abgabe 4.2}



\end{document}
